\begin{proof}
Use the canonical $s$-dimensional realization in the converse proof of Lemma~\ref{I:U:lem:compact}. Its observation operator is an isometry, so
\[
 \sum_{w\in\Sigma^*}T_w^*cc^*T_w=\UId_s.
\]
The identity $\UId_s=cc^*+\Psi(\UId_s)$ then yields
\[
 \Psi^N(\UId_s)=\UId_s-\sum_{|w|<N}T_w^*cc^*T_w\longrightarrow0
\]
in matrix norm. For a Hermitian matrix $X$, positivity and
$-\Unorm X\UId_s\preceq X\preceq\Unorm X\UId_s$ imply
$\Unorm{\Psi^N(X)}\le\Unorm X\,\Unorm{\Psi^N(\UId_s)}$.
Splitting a general complex matrix into Hermitian real and imaginary parts shows that the powers of $\Psi$ tend to zero as linear maps on the entire matrix space. In finite dimension this is equivalent to $\rho(\Psi)<1$.
For real realizations, apply the same argument to the complexification of the real matrix map.
The map $\Phi_T$ is the adjoint of $\Psi$ for the Hilbert--Schmidt inner product, and hence has the same spectral radius. Zero padding preserves strict stability; the zero operator has a zero-transition realization.
\end{proof}
